% Chapter 22: The Kadison--Singer Problem
\let\LocalMacro\relax

\chapter{Kadison--Singer Problem}

\section{The Problem}

\subsection{Informal}
In the 1950s, two mathematicians asked a question that sounds like a riddle: suppose you have a huge collection of vectors in space, each pointing in a different direction. Can you always split them into a few groups, each of which behaves almost as nicely as a set of mutually perpendicular axes? This question, which seemed purely abstract, turned out to have deep consequences for how we understand quantum measurements, signal processing, and even the structure of data.

\subsection{Formal}
\begin{conjecture}[Paving Conjecture]
For every $\epsilon > 0$, there exists an integer $r$ such that any orthogonal projection $P$ on $\mathbb{C}^n$ with diagonal entries $P_{ii} \leq \epsilon$ can be $r$-paved. That is, the index set $\{1, \ldots, n\}$ can be partitioned into $r$ subsets $S_1, \ldots, S_r$ such that for each $j$, the submatrix $P_{S_j}$ satisfies:
\begin{equation}
\|P_{S_j}\| \leq 1 - \eta
\end{equation}
for some $\eta > 0$ depending only on $\epsilon$.
\end{conjecture}

This conjecture implies that one can decompose a large projection into smaller parts whose operator norms are strictly bounded away from 1.

\subsection{Historical Context}
The Kadison--Singer problem remained unsolved for over 50 years. Its difficulty stemmed from its connection to diverse areas of mathematics, including C*-algebras, frame theory in signal processing, graph theory, and computer science.

\subsection{What Was Known}
The problem was posed in 1959 and was known to be equivalent to several other important conjectures in mathematics, including the Paving Conjecture in operator algebras and the Feichtinger conjecture in harmonic analysis. By the 2000s, partial results had been obtained using techniques from Banach space theory, but the full problem remained open for over 50 years.

\subsection{Why It Matters}
The Kadison--Singer problem sits at the crossroads of quantum physics, signal processing, and pure mathematics. A positive answer means that incomplete quantum measurements can always be completed consistently. This has implications for how we understand measurement in quantum mechanics and how we process incomplete data in engineering.

\subsection{Simplified Version}
The problem is equivalent to the Paving Conjecture: given an infinite-dimensional diagonal operator $D$ with bounded entries, can the basis vectors be partitioned into finitely many blocks such that each block is nearly orthogonal? This was verified numerically for many examples before the proof. Weaver had conjectured a finite-dimensional reformulation (Weaver's conjecture $KS_2$) that was already checked for small matrices. The breakthrough came from showing that Weaver's conjecture holds in general via the ``few moments'' method.

\section{Mathematical Theory}


\subsection{Prerequisites}
This chapter assumes familiarity with:
\begin{itemize}
\item Linear algebra (inner products, orthogonal projections, eigenvalues)
\item Operator theory and Hilbert spaces (bounded operators, pure states)
\item Basic graph theory (adjacency matrices, bipartite graphs)
\end{itemize}

\subsection{Frames and Discretization}
In signal processing, a \emph{frame} is a generalization of an orthonormal basis. A family of vectors $\{\phi_i\}_{i \in I}$ in a Hilbert space $H$ is a frame if there exist constants $A, B > 0$ such that:
\begin{equation}
A \|x\|^2 \leq \sum_{i \in I} |\langle x, \phi_i \rangle|^2 \leq B \|x\|^2
\end{equation}
for all $x \in H$. The Kadison--Singer problem is equivalent to the statement that any frame can be partitioned into a finite number of subsets, each of which is a frame with controlled bounds.

\subsection{Bipartite Ramanujan Graphs}
A $d$-regular graph is \emph{Ramanujan} if its adjacency matrix eigenvalues (excluding $\pm d$) have absolute value at most $2\sqrt{d-1}$. Ramanujan graphs are optimal expanders. In 2013, Marcus, Spielman, and Srivastava used the Kadison--Singer problem to prove the existence of bipartite Ramanujan graphs of all degrees.

\section{The Solution}

\subsection{The Big Idea}
Imagine you have a huge matrix and want to split it into two halves so that neither half is too ``large'' (in operator norm). You could try every possible split, but there are exponentially many---hopeless. The brilliant idea of Marcus, Spielman, and Srivastava was to *flip a coin for each row*: heads means put it in the first half, tails in the second. Then they studied the *average* behavior over all these random splits. They proved that the average is good enough that *some specific split* must also be good---using a clever argument about how the roots of polynomials behave when you combine them. It is like showing that if the average temperature in 1,000 cities is comfortable, then at least one city must be comfortable. The deep insight is that *existence can be proved without construction*: you don't need to find the good partition, just prove one exists.

\subsection{What Was Stopping Progress}
Two major families of techniques fell short. From operator algebras, one could study faces and projections in the state space of a C*-algebra, but these tools could not produce the combinatorial partitioning required by the Paving Conjecture. From signal processing, the Feichtinger conjecture---an equivalent reformulation about partitioning frames---could not be resolved using frame theory, Riesz basis arguments, or Banach space geometry. The difficulty was fundamental: all known approaches were either too qualitative (they could not quantify the norm bound) or too constructive (they required finding an explicit partition, which is an exponentially hard combinatorial problem).

\subsection{The Breakthrough}

Marcus, Spielman, and Srivastava's breakthrough was to \emph{stop trying to find a good partition} and instead prove one must exist. Their tool is \emph{interlacing families of polynomials}. Here is the core idea in a simple case: suppose you flip a fair coin for each of $m$ vectors, assigning each vector to group $+1$ or $-1$. The characteristic polynomial of the resulting matrix depends on the coin flips---it is a random polynomial. The \emph{expected} characteristic polynomial (averaged over all $2^m$ coin-flip sequences) turns out to have roots that interlace with a known polynomial. Interlacing means the roots alternate: if the known polynomial has roots $a_1 < a_2 < \cdots$, the expected polynomial has roots $b_1$ with $a_1 < b_1 < a_2$, and so on. This interlacing forces every individual coin-flip outcome to have its largest root bounded by the largest root of the known polynomial---so at least one specific signing works. The proof extends this from one signing to a recursive partition, using multivariate polynomials and barrier functions to track how roots move under rank-one updates.

\subsection{How It Works}

The Kadison--Singer problem was posed in 1959 and resisted solution for 54 years. The key reformulation was the Paving Conjecture: given an orthogonal projection $P$ on $\mathbb{C}^n$ with small diagonal entries, one must partition the index set so that each submatrix has operator norm strictly bounded away from 1. What was missing was a method to prove the \emph{existence} of such a partition without explicitly constructing it. Numerical evidence strongly supported the conjecture, and partial results had been obtained for special cases, but no general argument existed that could handle arbitrary projections.

Two major families of techniques were tried and fell short. From operator algebras, one could study faces and projections in the state space of a C*-algebra, but these tools could not produce the combinatorial partitioning required by the Paving Conjecture. From signal processing, the Feichtinger conjecture---an equivalent reformulation about partitioning frames---could not be resolved using frame theory, Riesz basis arguments, or Banach space geometry. The difficulty was fundamental: all known approaches were either too qualitative (they could not quantify the norm bound) or too constructive (they required finding an explicit partition, which is an exponentially hard combinatorial problem).

Marcus, Spielman, and Srivastava introduced \emph{interlacing families of polynomials}, a new algebraic framework. The core idea is: instead of trying to find the best partition directly, study the \emph{expected characteristic polynomial} of a random signing of the matrix. If one can show that the roots of this expected polynomial interlace with the roots of a known polynomial, then there must exist a specific signing whose characteristic polynomial has all roots bounded appropriately. This converts a combinatorial existence problem into a statement about polynomial roots---a far more tractable question. The proof introduced multivariate polynomials, barrier functions, and a recursive argument showing that rank-one updates preserve the interlacing property.

In 2013, Adam Marcus, Daniel Spielman, and Nikhil Srivastava solved the Kadison--Singer problem \cite{marcus2015a, marcus2015b}:

\begin{theorem}[Marcus--Spielman--Srivastava, 2013]
The Kadison--Singer Paving Conjecture is true. Specifically, for any vectors $v_1, \ldots, v_m$ in $\mathbb{C}^d$ such that $\sum_{i=1}^m v_i v_i^T = I$ and $\|v_i\|^2 \leq \epsilon$, there exists a partition of $\{1, \ldots, m\}$ into two sets $S_1, S_2$ such that:
\begin{equation}
\left\|\sum_{i \in S_j} v_i v_i^T\right\| \leq \frac{(1 + \sqrt{\epsilon})^2}{2}
\end{equation}
for $j=1,2$.
\end{theorem}

Their proof used a revolutionary new technique:

\begin{enumerate}
\item \textbf{Interlacing Families of Polynomials}: They showed that the roots of certain families of polynomials interlace, which guarantees the existence of a partition where the maximum root of the characteristic polynomial is small.
\item \textbf{Expected Characteristic Polynomials}: Instead of finding a partition directly, they studied the average characteristic polynomial over all random partitions, showing its roots are bounded by the roots of an associated univariate polynomial.
\item \textbf{Barrier Functions}: They used multivariate barrier functions to track and bound the roots of these polynomials under rank-one updates.
\end{enumerate}

\begin{highlightbox}
The Marcus--Spielman--Srivastava proof settled Kadison--Singer by converting an operator algebra problem into a question about the roots of polynomials. This ``method of interlacing families'' has since become a major tool in combinatorics and computer science.
\end{highlightbox}

\subsection{After the Proof}

The Kadison--Singer problem is completely settled. The Paving Conjecture is true, and all its equivalent formulations---the Feichtinger conjecture in harmonic analysis, the Weaver conjectures in combinatorics, the existence of certain Ramanujan graphs---follow as corollaries. There is no remaining doubt about the qualitative result.

What the proof unlocked, though, may be even more important than the result itself. The method of interlacing families of polynomials has become a major tool across combinatorics, theoretical computer science, and signal processing. It has been used to construct optimal expanders, analyze random structures, and prove existence results that were previously inaccessible. The interlacing families framework converts questions about eigenvalues of random matrices into questions about roots of polynomials, which are far more tractable. This connection between algebraic and probabilistic reasoning has opened up many new research directions.

Open questions remain at the quantitative level. The paving bound---how many pieces you need in the partition, and how tight the norm bound is---is not optimal. Current bounds are polynomial in $1/\epsilon$, but there is reason to hope for linear or near-linear bounds. Applications to quantum information theory, where the Kadison--Singer problem originally arose, are still being developed: the question of how to consistently complete partial quantum measurements has new formulations that the interlacing families method may help resolve.

\section{Impact}
\begin{itemize}
\item \textbf{Operator Algebras}: Resolved the 1959 Kadison--Singer C*-algebra conjecture.
\item \textbf{Signal Processing}: Proved that every frame can be paved, ensuring stable partitions in transmission codes.
\item \textbf{Computer Science}: Provided a deterministic construction of bipartite Ramanujan graphs of all degrees, solving a long-standing open problem in network design.
\end{itemize}

\section{Exercises}

\begin{exercise}[Basic]
Show that the union of two frames in $\mathbb{C}^d$ is a frame. If $A_1, B_1$ and $A_2, B_2$ are the frame bounds for the two frames, what are the bounds for the union?
\end{exercise}

\begin{exercise}[Intermediate]
Explain the relation between the Kadison--Singer problem and the construction of Ramanujan regular graphs.
\end{exercise}

\begin{exercise}[Challenge]
Prove that any orthogonal projection $P$ on $\mathbb{C}^n$ with diagonal entries $P_{ii} \leq \epsilon$ can be partitioned into two parts whose submatrices have norm bounded away from 1.
\end{exercise}

\begin{exercise}[Computational]
Write a Python/SageMath script to explore a concept from this chapter. See the appendix for starter code.
\end{exercise}

\section{Timeline}

\begin{tabularx}{\linewidth}{lX}
\toprule
\textbf{Year} & \textbf{Event} \\ \midrule
1959 & Kadison and Singer formulate the extension problem \cite{kadison1959} \\ 
1979 & Anderson reformulates the conjecture as the Paving Conjecture \cite{anderson1979} \\ 
1980s--2000s & Equivalent conjectures found in signal processing (Feichtinger conjecture) \\ 
2013 & Marcus, Spielman, and Srivastava resolve the conjecture using interlacing polynomials \cite{marcus2015a, marcus2015b} \\ 
2015 & Annals of Mathematics publishes the two-part proof \\ \bottomrule
\end{tabularx}

\section{Citations}

\begin{enumerate}
\item Kadison, R. V., \& Singer, I. M. (1959). ``Extensions of pure states.'' American Journal of Mathematics, 81(2), 383--400.
\item Anderson, J. (1979). ``Extreme points in secondary algebra.'' Journal of Functional Analysis.
\item Marcus, A., Spielman, D. A., \& Srivastava, N. (2015). ``Interlacing families I: Bipartite Ramanujan graphs of all degrees.'' Annals of Mathematics, 182(1), 307--325.
\item Marcus, A., Spielman, D. A., \& Srivastava, N. (2015). ``Interlacing families II: Mixed characteristic polynomials and the Kadison--Singer problem.'' Annals of Mathematics, 182(1), 327--350.
\end{enumerate}
